%%%%%%%%%%%%%%%%%%%%%%%%%%%%%%%%%%%%%%%%%
% Simple Sectioned Essay Template
% LaTeX Template
%
% This template has been downloaded from:
% http://www.latextemplates.com
%
%
%%%%%%%%%%%%%%%%%%%%%%%%%%%%%%%%%%%%%%%%%

%----------------------------------------------------------------------------------------
%	PACKAGES AND OTHER DOCUMENT CONFIGURATIONS
%----------------------------------------------------------------------------------------
%!TEX root = 00Main.tex

%----------------------------------------------------------------------------------------
%	PACKAGES AND OTHER DOCUMENT CONFIGURATIONS
%----------------------------------------------------------------------------------------
\PassOptionsToPackage{table}{xcolor}

\documentclass[12pt]{article} % Default font size is 12pt, it can be changed here

\usepackage[margin=1.2in,footskip=1cm]{geometry} % Required to change the page size to A4
\geometry{a4paper} % Set the page size to be A4 as opposed to the default US Letter
\usepackage{textcomp}
\usepackage{lipsum}
\usepackage{graphicx} % Required for including pictures
\usepackage{amsmath}
\DeclareMathOperator*{\argmin}{arg\,minimize}
\usepackage{courier}
\usepackage{bm} %boldsymbol mathmode
% \usepackage{nccmath} %Smaller fonts for big matrices
\usepackage{float} % Allows putting an [H] in \begin{figure} to specify the exact location of the figure
\usepackage{wrapfig} % Allows in-line images such as the example fish picture
% \usepackage{pgfgantt}
\usepackage{subcaption}
\usepackage{parskip} %Read in blank lines in text (allows for paragraphs)
\usepackage[super,numbers,sort&compress]{natbib}
\renewcommand*{\thefootnote}{(\arabic{footnote})} %Change footnote styling to avoid confusion with citations
\usepackage{array}
\usepackage{tikz}
\usepackage{xcolor}
\usepackage{listings} %Typesetting code
\usepackage{fancyhdr}	%Headers and footers
\usepackage{titlesec}	%Alter titles
\usepackage{sectsty}	
\usepackage{calc}
\usepackage{abstract}
\usepackage{appendix}
\usepackage[usenames,dvipsnames]{pstricks}
\usepackage{epsfig}
\usepackage{multirow}
\usepackage{longtable}
\usepackage[nottoc,numbib]{tocbibind}
\usepackage{multicol}
\usepackage{nicefrac}
\usepackage[greek.ancient,english]{babel}
\usepackage{gensymb}
\usepackage{mcode}
\usepackage{scrextend}
\usepackage{dsfont}
\usepackage{tikz}
\usepackage{booktabs}
\usepackage{hyperref}   % Hyperlinking
\usepackage{cleveref} % Has to be loaded last

%%%%%%%%%%%%%%%%%%%%%%%%%%%%  Defining colours  %%%%%%%%%%%%%%%%%%%%%%%%%%%%

\definecolor{BioEngTeal}{rgb}{0.43529411764,0.79607843137,0.82352941176}
\definecolor{BioEngGrey}{rgb}{0.56862745098,0.5725490196,0.58431372549}
\definecolor{mygreen}{RGB}{28,172,0} % color values Red, Green, Blue
\definecolor{mylilas}{RGB}{170,55,241}
\definecolor{RubineRed}{RGB}{206,0,88}

%%%%%%%%%%%%%%%%%%%%%%%%%%%%%%%%%%%%%%%%%%%%%%%%%%%%%%%%%%%%%%%%%%%%%%%%%%%%%

%%%%%%%%%%%%%%%%%%%%%%%%%%%  New environment for typesetting big matrices  %%%%%%%%%%%%%%%%%%%%%%%%

\newenvironment{mbmatrix}{\begin{medsize}\begin{bmatrix}}%
{\end{bmatrix}\end{medsize}}%

%%%%%%%%%%%%%%%%%%%%%%%%%%%%%%%%%%%%%%%%%%%%%%%%%%%%%%%%%%%%%%%%%%%%%%%%%%%%%%%%%%%%%%%%%%%%%%%%%%%%%

%%%%%%%%%%%%%%%%  Change colours of sections, subsections, figures and tables  %%%%%%%%%%%%%%%%%

\sectionfont{\color{BioEngTeal}}  % sets colour of sections
\subsectionfont{\color{BioEngGrey}}  % sets colour of subsections
\renewcommand{\abstractnamefont}{\color{BioEngTeal}\textbf}


\renewcommand{\bibname}{References}
\renewcommand{\thefigure}{\textcolor{BioEngTeal}{\bfseries\itshape\thesection.\arabic{figure}\,}}
\renewcommand{\thesubfigure}{\textcolor{BioEngTeal}{\bfseries\alph{subfigure}}}
\renewcommand{\figurename}{\textcolor{BioEngTeal}{\bfseries\itshape Figure}}
\renewcommand{\thetable}{\textcolor{BioEngTeal}{\bfseries\itshape\thesection.\arabic{table}}}
\renewcommand{\tablename}{\textcolor{BioEngTeal}{\bfseries\itshape Table}}
\renewcommand{\theequation}{\textcolor{BioEngTeal}{\bfseries\itshape\thesection.\arabic{equation}\,}}

\addto\captionsenglish{\renewcommand{\figurename}{\textcolor{BioEngTeal}{\bfseries\itshape Figure}}}
\addto\captionsenglish{\renewcommand{\tablename}{\textcolor{BioEngTeal}{\bfseries\itshape Table}}}
%%%%%%%%%%%%%%%%%%%%%%%%%%%%%%%%%%%%%%%%%%%%%%%%%%%%%%%%%%%%%%%%%%%%%%%%%%%%%%%%%%%%%%%%%%%%%%%%%




%%%%%%%%%%%%%%%%%%%%%%%%%%%%%%% Define macro for C++ symbol  %%%%%%%%%%%%%%%%%%%%%%%%%%%%%%%%%%%%

\newcommand{\CC}{C\nolinebreak\hspace{-.05em}\raisebox{.4ex}{\bf +}\nolinebreak\hspace{-.10em}\raisebox{.4ex}{\bf +}\,\,}
\def\CC{{C\nolinebreak[4]\hspace{-.05em}\raisebox{.4ex}{\bf ++}}\,\,}

%%%%%%%%%%%%%%%%%%%%%%%%%%%%%%%%%%%%%%%%%%%%%%%%%%%%%%%%%%%%%%%%%%%%%%%%%%%%%%%%%%%%%%%%%%%%%%%%%


%%%%%%%%%%%%%%%%%%%%%%%%%%  Defining norm symbol  %%%%%%%%%%%%%%%%%%%%%%

\newcommand{\norm}[1]{\left\lVert#1\right\rVert}

%%%%%%%%%%%%%%%%%%%%%%%%%%%%%%%%%%%%%%%%%%%%%%%%%%%%%%%%%%%%%%%%%%%%%%%%


%%%%%%%%%%%%%%%%%%%%%%%%%%%%%%%%%  Defining footer  %%%%%%%%%%%%%%%%%%%%%%%%%%%%%%%%%%

\pagestyle{fancy}
\fancyhf{}
\renewcommand{\footrulewidth}{0.5pt}
\renewcommand{\headrulewidth}{0pt}

\fancyfoot[R]{\color{BioEngTeal}\textbf{\thepage}}
\fancyfoot[L]{\raisebox{-.75\height}{\includegraphics[scale=1]{Pictures/bioengineering_logo_left.eps}}}

%%%%%%%%%%%%%%%%%%%%%%%%%%%%%%%%%%%%%%%%%%%%%%%%%%%%%%%%%%%%%%%%%%%%%%%%%%%%%%%%%%%%%%





%%%%%%%%%%%%%%%%%%%%%%%  Change font to Helvetica %%%%%%%%%%%%%%%%%%%%%%%%%%%

% % Include the helvet package
% \usepackage{helvet}

% % Set the default font to be sans-serif
% \renewcommand*{\familydefault}{\sfdefault}
% \makeatletter
% \newcommand{\thickhline}{%
%     \noalign {\ifnum 0=`}\fi \hrule height 2pt
%     \futurelet \reserved@a \@xhline
% }
% \newcolumntype{"}{@{\hskip\tabcolsep\vrule width 1pt\hskip\tabcolsep}}
% \makeatother
%%%%%%%%%%%%%%%%%%%%%%%%%%%%%%%%%%%%%%%%%%%%%%%%%%%%%%%%%%%%%%%%%%%%%%%%%%%%%%%%%%%%%

%%%%%%%%%%%%%%%%%%%%%%%  Fonts  %%%%%%%%%%%%%%%%%%%%%%%%%%%
\usepackage[scale=.75]{sourcecodepro} % monospace font
\usepackage{courier}
\renewcommand{\ttfamily}{pcr}
\usepackage{gfsbaskerville}
\usepackage[osf]{Baskervaldx} % tosf in text, tlf in math
\usepackage[T1]{fontenc}
% \usepackage[baskervaldx,cmintegrals,bigdelims,vvarbb]{newtxmath} % math italic letters from Baskervaldx
% \usepackage[cal=boondoxo]{mathalfa} % mathcal from STIX, unslanted a bit
% \usepackage{ebgaramond-maths}
%%%%%%%%%%%%%%%%%%%%%%%%%%%%%%%%%%%%%%%%%%%%%%%%%%%%%%%%%%%%%%%%%%%%%%%%%%%%%%%%%%%%%

\linespread{1.2} % Line spacinggit gitg

\setlength\parindent{0pt} % Uncomment to remove all indentation from paragraphs

\graphicspath{{Pictures/}} % Specifies the directory where pictures are stored

%%%%%%%%%%%%%%%%%%%%%%%%%%%%%%%%%%%%%%%%%%%%%%%%%%%%%%%%%%%%%%%%%%%%%%%%%%%%%%%%%%%%%%



%%%%%%%%%%%%%%%%%%%%%%%%%%%%%  Typesetting code  %%%%%%%%%%%%%%%%%%%%%%%%%%%%%%%%%%%%%
\definecolor{listinggray}{gray}{0.9}
\definecolor{lbcolor}{rgb}{0.9,0.9,0.9}
\definecolor{Darkgreen}{rgb}{0.13,0.545,0.13}

\usepackage[numbered,framed]{matlab-prettifier}

\usepackage{filecontents}

\let\ph\mlplaceholder % shorter macro
\lstMakeShortInline"

\lstset{
  style              = Matlab-editor,
  basicstyle         = \mlttfamily,
  escapechar         = ",
  mlshowsectionrules = true,
}


% Command for placeholder
\newcommand\placeholder[1]%
{%
    \bgroup
        \normalfont\upshape\color{RubineRed}%
        < {\itshape #1\/} >%
    \egroup
}


% Command for error
\newcommand\error[1]%
{%
    \bgroup
        \color{red}%
        #1\/
    \egroup
}



\lstset
{%
    escapechar=`,
}
%%%%%%%%%%%%%%%%%%%%%%%%%%%%%%%%%%%%%%%%%%%%%%%%%%%%%%%%%%%%%%%%%%%%%%%%%%%%%%%%%%%%%%%%
\begin{document}
\hrule

\begin{figure}[t]
	\begin{subfigure}[b]{0.40\linewidth} 					
    \includegraphics[height=0.8cm]{logo.jpg}     
  \end{subfigure}
    \hfill
	\begin{subfigure}[B]{0.40\linewidth} 						
    \includegraphics[height=0.8cm]{Pictures/bioengineering_logo_right.eps}
  \end{subfigure}
    \hfill
\end{figure}

\vspace{1cm}


\begin{center}
{\large \textbf{Mini-project 2018: Image processing}} % Title of your document
\\
\textsc{Solutions}
\end{center}

\section{Rotation}
\begin{minipage}{\textwidth}
\begin{lstlisting}
function output_image = rotate_image(input_image, theta)
   
    % Define the tranformation matrix
    theTransform = inv([cos(theta) sin(theta);...
                       -sin(theta) cos(theta)]); 
    
    % Get size of input image
    [ny, nx] = size(input_image); 
    
    % Determine centre of input image
    xCent = round(nx/2); yCent = round(ny/2); 
    
    % Preallocate output image
    output_image = zeros(ny, nx); 
    
    % Preallocate source pixel
    src = zeros(2, 1); 
    
    % Loop through each pixel of output image and find
    % coordinates of input image for source pixel
    for ix = 1: nx 
        for iy = 1: ny 
            src=round(theTransform*([ix; iy]-[xCent; yCent])+...
                                             [xCent; yCent]); 
            xSrc = src(1); 
            ySrc = src(2);
        
            % Check if source coordinates are valid
            if (~(xSrc<1 || ySrc<1 || xSrc>nx || ySrc>ny))
                output_image(iy, ix) = input_image(ySrc, xSrc); 
            end
        end
    end
end
\end{lstlisting}
\end{minipage}

\section{Edge detection}

\subsection{Loops}
\begin{minipage}{\textwidth}
\begin{lstlisting}
function output_image = detect_edges(input_image)
    
    % Get size of input image
    [ny, nx] = size(input_image);
    
    % Preallocate output image
    output_image = zeros(ny-1,nx-1);
    
    % Loop though each pixel of the output image and apply the edge
    % detection formula to the input image
    for iy = 1:ny-1
        for ix = 1:nx-1
            output_image(iy,ix) = (abs(input_image(iy,ix) -... 
                                       input_image(iy+1,ix)) + ...
                                   abs(input_image(iy,ix) -... 
                                       input_image(iy,ix+1)))/2;
        end
    end 
end
\end{lstlisting}
\end{minipage}

\subsection{Vectorised}
\begin{minipage}{\textwidth}
\begin{lstlisting}
function output_image = detect_edges_vectorised(input_image) 

    % get size of input image
    [ny, nx] = size(input_image); 

    % Create vectors used to index input image
    yv = 1:ny-1; 
    xv = 1:nx-1; 

    % Apply edge detection formula
    output_image=...
    			(abs(input_image(yv, xv) -input_image(yv, xv+1)) +... 
                 abs(input_image(yv, xv)-input_image(yv+1, xv)))/2; 

end
\end{lstlisting}
\end{minipage}
\newpage
\section{Blurring}
\subsection{Loops}
\begin{minipage}{\textwidth}
\begin{lstlisting}
function output_image = blur_image(input_image)

    % get size of input image
    [ny, nx] = size(input_image);
    
    % Preallocate output image
    output_image = zeros(ny-2,nx-2);
    
    % Loop through each pixel of the output image and apply the 
    % blurring formula to the input image (by looping through 
    % neighbouring pixels)
    for iy = 1:ny-2
        for ix = 1:nx-2
            for i = 0:2
                for j = 0:2
                    output_image(iy,ix) = output_image(iy,ix) +... 
                                          input_image(iy+i,ix+j)/9;
                end
            end
        end
    end
end
\end{lstlisting}
\end{minipage}

\subsection{Vectorised}
\begin{minipage}{\textwidth}
\begin{lstlisting}
function output_image = blur_image_vectorised(input_image) 

    % Ensure input is double precision, if it is uint
    % pixel values will clip at 255
    input_image = double(input_image);

    % get size of input image
    [ny, nx] = size(input_image); 

    % Create vectors used to index input image
    yv = 1:ny-2; 
    xv = 1:nx-2; 

    % Apply edge blurring formula
    output_image=...
        (input_image(yv, xv) + input_image(yv+1, xv) +... 
         input_image(yv + 2, xv) + input_image(yv, xv+1) +... 
         input_image(yv, xv+2) + input_image(yv+1, xv+1) + ...
         input_image(yv+1, xv+2) + input_image(yv+2, xv+1) +...
         input_image(yv+2, xv+2))/9;

end
\end{lstlisting}
\end{minipage}




\end{document}